\documentclass[11pt]{article}
\usepackage[T1]{fontenc}
\usepackage[utf8]{inputenc}
\usepackage{lmodern}
\usepackage[margin=0.9in]{geometry}
\usepackage{amsmath,amssymb,amsthm,mathtools}
\usepackage{booktabs,longtable,tabularx,array}
\usepackage{multicol}
\usepackage{enumitem,xcolor,microtype}
\usepackage{listings}
\usepackage[colorlinks=true,linkcolor=blue!45!black,urlcolor=blue!45!black,citecolor=blue!45!black]{hyperref}
\usepackage{fancyhdr}
\pagestyle{fancy}
\fancyhf{}
\fancyhead[L]{\small Further Lean proof simplifications}
\fancyhead[R]{\small ProveIt / 26 October 2026}
\fancyfoot[C]{\thepage}
\setlength{\headheight}{14pt}
\setlength{\emergencystretch}{3em}
\setlist[itemize]{topsep=5pt,itemsep=3pt}
\setlist[enumerate]{topsep=5pt,itemsep=3pt}
\newcommand{\code}[1]{\texttt{\detokenize{#1}}}
\newcommand{\N}{\mathbb N}
\newcommand{\Z}{\mathbb Z}
\newcommand{\zm}[1]{\Z/#1\Z}
\newcommand{\repo}{\url{https://github.com/VladimirReshetnikov/ProveIt}}
\newcommand{\pin}{\texttt{e1c8914bc7064d9a8e7c06171a18069a949badd0}}
\newcommand{\leanpath}{\texttt{Computability/DiophantineConsistency/Lean/}}
\newtheorem{proposition}{Proposition}[section]
\newtheorem{lemma}[proposition]{Lemma}
\newtheorem{corollary}[proposition]{Corollary}
\lstset{basicstyle=\small\ttfamily,columns=fullflexible,keepspaces=true,
  breaklines=true,frame=single,rulecolor=\color{black!18},
  backgroundcolor=\color{black!2},aboveskip=9pt,belowskip=9pt,
  showstringspaces=false}
\title{Four further directions for shorter Lean proofs\\[5pt]
  \large Reflection, fixed prime support,\newline
  translated subset sums, and prime-power order}
\author{Research and verification packet for Vladimir Reshetnikov}
\date{26 October 2026}
\begin{document}
\maketitle
\begin{abstract}
This packet examines the pinned ProveIt source after treating every proposal
in the requested research corpus as already applied. The initial review of
23 placed packets and five incoming archives was extended to all 36 packets
at the late audit cutoff. It retains four additional, narrowly described directions. The most
local one removes an unnecessary upper-field hypothesis from reflected-source
normalization. The other directions transport controller coprimality through a
fixed set of prime factors, realize prescribed finite-subset sums by
translation, and characterize all prime-power divisibility before extracting
multiplicative order. The accompanying Lean files check mathematical
interfaces; the validation section distinguishes these from a rebuilt
production module. No general source-line reduction or whole-project build
is claimed merely because a client theorem becomes shorter.
\end{abstract}

\section{Scope, evidence, and the counterfactual baseline}

The source pin is \pin. Its commit timestamp is 26 October 2026,
18:30:35 UTC. Primary source links identify this immutable snapshot; the late comparison
below uses its own explicitly recorded head. The requested Lean subtree contains 690 \code{.lean} files,
including check files, with 86,351 physical lines. These inventory figures
describe the entire materialized subtree, not the amount of code a proposed
refactor has removed. The toolchain is Lean 4.32.0; the exact Mathlib revision
is \code{81a5d257c8e410db227a6665ed08f64fea08e997}.

The novelty comparison includes placed packets \code{2026-10-26a} through
\code{2026-10-26w} and all five ZIPs in \code{docs/incoming} at the pin.
The report inventory contains 82 TeX input files, totaling 1,086,021 bytes;
the five ZIPs contain 1,042 nondirectory members. Several inputs and members
are duplicates, supporting material, or old source snapshots. These are not
counts of distinct ideas. The audit read recommendation sections, delivery
notes, exclusion ledgers, and candidate-specific supporting prose and code.
It did not execute programs supplied by prior packets, and does not claim
that every duplicated snapshot was reviewed line by line. The exact inventory
and a conservative exclusion ledger accompany this report.\cite{prior,incoming}

\paragraph{Late refresh.}
The final exclusion cutoff is
\code{ae736b383e59c32c3c3f81910ef1c1197cbfbd44}.
At that head the research directory contains 36 packets, \code{a} through
\code{zj}. The five previously audited incoming archives are now placed as
\code{x} through \code{zb}; eight additional packets, \code{zc} through
\code{zj}, were screened separately. Their complete recursive tree inventories
were retrieved. The audit fetched 71 selected report, note, review, and
candidate-code files, totaling 575,496 bytes, plus both directory indexes;
all 73 fetched texts match their Git blob hashes. The incoming directory
contains only its README at this cutoff. The eight direct source targets
checked at both heads are byte-identical; the new safety, controller-subset,
position, rephasing, and population modules were inspected separately.
The precise bounded review and inventory are in
\code{evidence/late-refresh.md} and
\code{evidence/late-corpus-manifest.json}.\cite{lateprior,lateincoming}

\paragraph{What ``additional'' means here.}
An earlier proposal remains excluded if it is pending, uncompiled, optional,
or recorded only in a further-research section. A new file name, a new tactic,
or a new successful compilation does not make the old idea additional.
The claim is a bounded comparison with the specified repository corpus;
it is not a claim of global mathematical originality.
The assumed-applied baseline governs which mechanisms are additional.
Source counts and patch measurements refer to the actual pinned files;
they are not additive savings after a hypothetical composition of every
earlier patch, some of which offer alternative interfaces.

\paragraph{Important exclusions.}
This review does not propose scanner iteration helpers, generic PR closure
automation, shared circuit tails, exact carry quotient/remainder interfaces,
fixed-word normalization uniqueness, standard aligned-cell equivalences,
cyclic packing homomorphisms, or stronger-to-weaker modulus projection again.
It also excludes the newer helical-coordinate, row-recognition,
signal-statistics, affine-rank, and signed-prefix ideas in the incoming ZIPs.
In particular, computing the kernel once modulo a product of three primes
was considered and rejected as a headline: it is too close to the earlier
modulus-projection proposal.

The full cell/field-coordinate and numeric control-update direction was
also dropped. Packet \code{w} already gives the more general band/time
equivalence, and packet \code{zd} supplies the independent control-slot
overwrite representation and its invariants. The terminal-band convolution
consolidation in packet \code{zh} is likewise excluded. Those decisions
are reflected in the four retained mechanisms below.\cite{updates,terminal}

\begin{longtable}{p{0.03\textwidth}p{0.24\textwidth}p{0.37\textwidth}p{0.26\textwidth}}
\toprule
 & Additional obligation & Main current clients & Intended simplification\\
\midrule
\endhead
A & Normalize reflection at every target field & \code{CurrentReflectedGeometry2011}, \code{CurrentWordTopBounds2011} & One source formula handles low, top, and multi-period targets.\\
B & Keep unit status fixed while a modulus gains prime powers & \code{CurrentControllerUnits2011} & Separate one fixed unit certificate from variable exponents.\\
C & Translate any finite seed to a desired sum & \code{ControllerIntervalSums}, \code{ControllerOrbitTransport} & Existence clients no longer depend on triangular-sum evaluation.\\
D & Characterize divisibility at arbitrary exponents & \code{PrimePowerOrder} & Derive order through a single universal characterization.\\
\bottomrule
\end{longtable}

\section{A: remove the unnecessary target bound in reflection}

\subsection{The current proof boundary}

The theorem \code{CurrentReflectedGeometry2011.reflected_source} currently
assumes \(h\leq1004\) and \(p\leq985\). For a base \(b\) aligned modulo 2011
and a word \(c\) supported in the low fields, it proves
\begin{equation}\label{eq:reflection}
c\bigl(b+h-(1005-p)\bigr)=
\begin{cases}
c\bigl(b+(h+p-1005)\bigr),&1005\leq h+p,\\
0,&h+p<1005.
\end{cases}
\end{equation}
Positions on the two sides live in \(\operatorname{ZMod}(2011N)\);
the subtractions displayed inside natural casts are natural subtraction.
The bounds are therefore part of the proof, not an invitation to treat
truncated subtraction as integer subtraction.\cite{reflection}

The upper bound on \(h\) is unnecessary. The proof already splits on
\(1005\leq h+p\). In the positive branch, the cast identity is valid for
arbitrarily large \(h\). In the negative branch, the branch hypothesis itself
implies \(h<1005\), which is stronger than the only needed local bound.
The coefficient bound can also be weakened to \(p\leq1005\).

\begin{proposition}\label{prop:allfields}
Let \(N>0\). Suppose \(b\in\operatorname{ZMod}(2011N)\) has field residue zero
and \(c(k)=0\) whenever \(1004<\operatorname{val}(k)\bmod2011\).
Then \eqref{eq:reflection} holds for every \(h\in\N\) and every \(p\leq1005\).
\end{proposition}
\begin{proof}
If \(h+p\geq1005\), both natural subtractions are justified by their stated
inequalities, and the equality is ring arithmetic after casting.
Otherwise put \(d=1005-p\). Then \(h<d\leq1005\), so the field residue of
\(b+h-d\) lies between 1006 and 2010. It is outside the permitted support,
and the source digit is zero. This reasoning needs no global upper bound
on \(h\).
\end{proof}

\subsection{How this changes the existing proofs}

First promote the source formula to \code{reflected_source_all} and lift it
through the existing finite sum to \code{reflected_cell_all}. The old theorem
types can initially be retained as wrappers. This provides a reversible
migration path while making the stronger result available to new clients.

Late packet \code{2026-10-26zh} already consolidates the exact terminal-band
response, including the 2009 and 2010 endpoints in
\code{CurrentWordTopBounds2011}. That consolidation is assumed
applied and is not claimed here. Its high-field coordinate helper still
requires a high-field hypothesis. The remaining change is the all-offset
guard formula, including arbitrary targets through multiple periods,
with the weakened input-position bound and compatibility wrapper. Any
further top-client migration must be measured against the terminal-band
refactor, not against its older separate endpoint proofs.\cite{topbounds,terminal}

This does not say that every downstream low-field theorem generalizes.
For example, \code{source_le}, \code{source_nonselector_coefficient}, and
the low-band coefficient budget in \code{reflected_lt} still use their
upper bounds for real mathematical reasons. Generalizing those by deleting
the same hypothesis would be unsound. The change is confined to source
normalization and its finite-sum lifting.

\paragraph{A useful guard test.}
The weakened coefficient bound is meaningful: permitting \(p>1005\) is
generally false because \(1005-p\) truncates to zero. For \(p=1006,h=0\),
the original source is at the base while the proposed positive-branch
normal form is one field later. An arbitrary supported word need not have
equal digits there. The prototype includes this boundary distinction.

\paragraph{Novelty and payoff.}
Earlier packets centralize field residues and carry projections. Even after
those are assumed available, the artificial restriction to low target fields
remains. The additional result is the removal of that restriction and a
single guarded source identity over all offsets. It is a small interface
strengthening, not a prediction of a large line-count reduction. Retaining
all compatibility wrappers increases source initially; deleting the
redundant top-only path requires an explicit API decision and fresh client
checks.

\section{B: separate fixed prime support from varying exponents}

\subsection{The repeated controller obligation}

The actual controller uses
\[
P=120199,\qquad d=47\cdot2011,\qquad
q_n=B^{P^n},\qquad A_r=K+B^{P^r},
\]
and the coefficient
\begin{equation}\label{eq:gamma}
\Gamma_{r,n}=\rho^{977}(q_n^2-1)(1+q_nA_r).
\end{equation}
The current proof establishes separate residues and nonzero tests at 47,
2011, and \(P\), then combines the resulting coprimality statements for the
modulus \(dP^n\). The width exponent \(r\) and cell exponent \(n\) are
independent; both can be zero.\cite{units}

The unit question only depends on the prime divisors of the modulus, not
on their multiplicities. The fixed modulus \(R=dP\) therefore suffices to
control every \(dP^n\). A useful reusable lemma is even more general:

\begin{proposition}\label{prop:thicken}
If \(x\equiv b\pmod{dp}\) and \(\gcd(b,dp)=1\), then
\(\gcd(x,dp^n)=1\) for every natural \(n\).
No primality or squarefreeness assumption on \(p\) is required.
\end{proposition}
\begin{proof}
Congruence preserves being a unit modulo \(dp\), so \(x\) is coprime to
both \(d\) and \(p\). Coprimality with \(p\) implies coprimality with
\(p^n\), and multiplying the two coprime factors gives the conclusion.
\end{proof}

The direction in this statement matters. When \(n=0\), the target modulus
is only \(d\); the converse need not hold. Replacing the implication by an
unqualified equivalence would introduce a boundary-case error.

\subsection{Freeze the coefficient, not the independent dimensions}

The existing power reductions show, at each of the three fixed primes,
\(q_n\equiv B\) and \(B^{P^r}\equiv B\). Combining these congruences at the
fixed product modulus gives
\[
\Gamma_{r,n}\equiv
\Gamma_*:=\rho^{977}(B^2-1)(1+B(K+B))\pmod R.
\]
Thus the final unit theorem can have one fixed arithmetic certificate
\(\gcd(\Gamma_*,R)=1\), one parameterized polynomial congruence, and one
application of Proposition~\ref{prop:thicken}. The existing checked kernel
and power residues remain the evidence for the fixed certificate and
congruences. They are not replaced by an unproved coprimality assumption. To avoid a
circular proof, build the fixed certificate directly from fixed residue
facts. Do not replace the variable coprimality endpoints with the new
radical argument and then use those same endpoints to justify its base
certificate.

The supplied \code{rankFactor_cast_congr} theorem retains explicit guards
\(1\leq q^2\) and \(1\leq B^2\). They justify casting the natural
subtractions in \eqref{eq:gamma}. The variables standing for \(q_n\) and
\(B^{P^r}\) remain distinct throughout the theorem.

\paragraph{Why this goes beyond divisor projection.}
Earlier packets already recommend deriving a weaker congruence from a
stronger modulus. Here the conclusion concerns \(dp^n\), which generally
does not divide \(dp\). The bridge is invariance of unit status under
raising the exponents of existing prime factors. This is the additional
mathematical step. Simply fusing three kernel computations into one modulus
without this step was excluded during the audit.

\paragraph{Expected proof reorganization.}
Use the fixed-support certificate as the main path to
\code{coprime_Gamma_width_cells} and \code{isUnit_Gamma_width_cells}.
Retain individual residue theorems if they have independent clients or
form part of the public diagnostic API. If they are kept, their source
cannot be counted as deleted. The generic checked theorem does not by
itself certify the concrete current \(\Gamma_*\) or the complete integration.

\section{C: translate a finite seed instead of evaluating its triangular sum}

\subsection{The sum changes predictably under translation}

Let \(G\) be an additive commutative group, let \(S\subseteq G\) be finite,
and write \(S+t=\{x+t:x\in S\}\). Translation is injective, so
\begin{equation}\label{eq:translate}
|S+t|=|S|,\qquad
\sum_{x\in S+t}x=\sum_{x\in S}x+|S|\,t.
\end{equation}
Therefore, if multiplication by \(|S|\) is surjective on \(G\), every
desired sum is realized by a translate of the same seed.

\begin{proposition}\label{prop:seed}
Let \(S\) be a finite subset of \(\operatorname{ZMod}T\) with cardinality
\(a\), and assume \(\gcd(a,T)=1\). For every target \(u\) there is a
translate \(S+t\) of cardinality \(a\) and sum \(u\).
\end{proposition}
\begin{proof}
The residue of \(a\) is a unit. Take
\(t=a^{-1}(u-\sum S)\) and apply \eqref{eq:translate}.
\end{proof}

If only the existence of some \(a\)-element subset is required and
\(a\leq T\), choose any such subset from the finite universe and then
translate it. There is no need to express its sum as \(a(a-1)/2\), or even
to choose consecutive coordinates.

\subsection{The precise benefit for the current source}

\code{ControllerIntervalSums} currently constructs an interval, proves its
injectivity and cardinality, evaluates its triangular sum in \(\N\), defines
\code{startFor}, and then proves the prescribed-sum result. This is a sound
explicit construction. The proposed alternate path is for
\code{exists_subset_card_sum} and downstream clients that consume only
cardinality and sum.\cite{intervals,orbit}

The distinction is essential for honest accounting. The existing
\code{exists_subset_card_sum} theorem already has a one-line proof using
the explicit construction. Replacing that line by a longer standalone
proof is not a local source shortening. The advantage is that an
existence-only dependency layer can avoid the interval-specific arithmetic,
and can accept arbitrary seeds when future constraints make another shape
more convenient. If the public \code{startFor} and \code{targetInterval}
formulas must remain, their triangular-sum proof remains too.

The Lean prototype proves the translation statement first over an arbitrary
additive commutative group, then derives the \code{ZMod} instance. It tests
an even modulus, a non-interval seed, and the zero-ring boundary \(T=1\).
It never divides by two in \(\operatorname{ZMod}T\). A unit inverse is used
only for the coefficient \(a\), whose coprimality premise is explicit.

\paragraph{A limit relevant to future controller work.}
Translation preserves size and controls the sum. It does not automatically
preserve membership in an arbitrary allowed-position set, safe mutation,
absence of carries, or a nonwrapping inequality. Such conditions require
separate invariance or safety lemmas. The result should not be presented as
a substitute for the separate safety argument. At the later audited head,
\code{CurrentDummySafety2011} already proves safe nonwrapping and exact
normalized mutation under explicit placement bounds. The present
translation theorem neither re-proves nor removes those premises.\cite{safety}

\section{D: prove arbitrary-exponent divisibility before multiplicative order}

\subsection{The stronger intermediate statement}

\code{PrimePowerOrder.exact_lift} currently specializes the lifting-the-
exponent calculation to exponents of the form \(p^k\), and
\code{order_succ} then separates the zero and successor cases of its index.
A more reusable normal form keeps the exponent arbitrary.\cite{order}

\begin{proposition}\label{prop:lte}
Let \(p\) be an odd prime, suppose \(p\mid g-1\), and suppose
\(p^2\nmid g-1\). For all \(n,k\in\N\),
\begin{equation}\label{eq:lteiff}
p^{n+1}\mid g^k-1\quad\Longleftrightarrow\quad p^n\mid k.
\end{equation}
\end{proposition}
\begin{proof}
The assumptions force \(g>1\) and \(p\nmid g\).
The standard odd-prime LTE identity gives
\(v_p(g^k-1)=1+v_p(k)\). Express both divisibilities through extended
multiplicity and cancel the finite summand one. For \(k=0\), both
divisibilities hold and extended multiplicity is infinite, so the same
formulation covers the boundary. Natural subtraction is justified by
\(g\geq1\), rather than silently replaced by integer subtraction.
\end{proof}

The order statement becomes a characterization of exactly the exponents
that annihilate the group element:
\[
(g\bmod p^{n+1})^k=1
\quad\Longleftrightarrow\quad p^n\mid k.
\]
Using \code{orderOf_dvd_iff_pow_eq_one}, divisibility in both directions
identifies the order as \(p^n\). The least-exponent argument no longer
requires a separate zero/successor branch.

\subsection{Compatibility and accounting}

The existing \code{exact_lift} theorem can be recovered by specializing
\eqref{eq:lteiff} to \(k=p^j\), while the public
\code{order_prime_power} endpoint can retain its original type. The late corpus already proposes a finite prime-factor certificate for
exact multiplicative order and a Mersenne exponent classification. Those
are treated as baseline. The new content here is the universal LTE
divisibility equivalence for arbitrary exponents, not another finite order
certificate. The value
of the new bridge is broader reuse of the full divisibility statement;
it is not enough to count only the shortened \code{order_succ} body and
ignore the new theorem or compatibility wrappers. The accompanying patch is a bounded production experiment for this
direction: the complete original and replacement source, all existing
module checks, added exponent controls, and public axiom comparisons were
replayed successfully. Standalone kernels have a different status.\cite{updates,mersenne}

The primary patch preserves the original cast-normalization proof verbatim
to avoid claiming routine cast cleanup as a new idea. It changes the full
file from 80 to 77 physical lines, including the new universal lemma.
UTF-8 size increases from 3,415 to 3,425 bytes. Thus the measured result is
three fewer lines and a stronger reusable interface, not byte compression.
All six original public theorem statements and their axiom sets are preserved.

Oddness and exact first-order lifting are indispensable. For example,
the familiar behavior of units modulo powers of two differs, and a base
already congruent to one modulo \(p^2\) has a smaller order than this
formula predicts. No such hypotheses are weakened by the proposed refactor.

\section{Validation, source accounting, and adoption order}

The machine-readable receipt in \code{validation/results.json} is
authoritative for exact file hashes, compiler invocations, exit codes, and
the distinction between prototype and production checks. The portable
replay instructions are in \code{README.txt}. All files labelled checked
must be checked with the repository-pinned Lean and Mathlib, not merely
accepted by a Python calculation or a source parser.

\begin{longtable}{p{0.20\textwidth}p{0.33\textwidth}p{0.38\textwidth}}
\toprule
Artifact & What was checked & What the check does not establish\\
\midrule
\endhead
Reflection kernel & Pinned Lean accepted the all-offset theorem, sparse-support adapter, sum lifting, wrappers and boundary controls. & The compatibility-preserving production patch was source-reviewed; the complete imported project module and its clients were not rebuilt.\\
Prime-support kernel & Pinned Lean accepted fixed-residue unit transfer and the guarded polynomial congruence, including exponent zero and a composite thickened factor. & It does not evaluate the literal current kernel or discharge the concrete fixed unit certificate.\\
Subset-translation kernel & Pinned Lean accepted the additive-group theorem, the modular unit instance, and the existence endpoint. & It does not replace the public exact interval definitions or establish safe physical mutation.\\
Prime-power replacement & Complete original/replacement source, existing checks, added arbitrary/zero-exponent controls, and axiom comparisons all passed. & This is a focused module replay, not a whole-repository build or a downstream auxiliary-order client rebuild.\\
\bottomrule
\end{longtable}

\noindent The prime-power patch preserves all six original public statement
texts and axiom sets, and adds \code{lift_divisibility}. Including the new
helper, the full file changes from 80 to 77 physical lines and from 3,415
to 3,425 bytes. The reflection proposal preserves the old source theorem
as a wrapper and changes its full module from 350 to 351 lines; its shorter
core proof is not advertised as a net source deletion. No net deletion is
claimed for the two larger architectural prototypes. The logged compiler
runs contain no \code{sorryAx}; standard Lean/Mathlib axioms are recorded
in the receipts. Linter diagnostics in concrete reflection examples do not
affect kernel acceptance.


\paragraph{Axioms and executable behavior.}
The audit records \code{#print axioms} for the new principal theorems.
Standard Lean/Mathlib axioms such as propositional extensionality,
classical choice, and quotient soundness are distinguished from an
unproved project assumption. The bundle does not use \code{sorry},
\code{admit}, or a new axiom as a substitute for a proof. Any claim that
a production patch preserves public theorem types is supported separately
by the patch and signature comparison.

\paragraph{Recommended adoption sequence.}
Start with A, because it identifies an unnecessary hypothesis at a precise
existing interface. Next assess D as a standalone Mathlib-only module
refactor with a complete checked theorem boundary.
B merits a separate controller integration because the concrete fixed
unit certificate and the product-modulus congruences must be wired to
existing facts. C is best introduced when splitting an existence-only
dependency layer or accepting a seed constrained by another argument.

For every migration, charge new helper definitions, wrapper theorems,
new imports, and retained compatibility code. Measure source reduction
on the final touched production files. Compiler time, peak memory, and
compiled-object size require their own measurements; none follows from
source length. A stronger and clearer interface may be worthwhile even
when its first adoption grows the file. The production file targeted by
the prime-power patch is byte-identical at both audited heads.

\section{Further research questions}

\begin{enumerate}
\item \textbf{Which field bounds are genuinely geometric?} After A,
classify each use of a low-target bound as either source-coordinate
normalization, support exclusion, or coefficient budgeting. Can the
remaining top-bound proofs be obtained from a single reflected support
description without reintroducing the already-proposed generic support
framework as a new result?
\item \textbf{Can fixed prime support become an explicit controller API?}
Package the fixed unit certificate, the two independent power reductions,
and the subtraction guards so that future changes in the exponents do not
reopen the kernel-residue proof. Check whether existing individual residue
endpoints are needed outside diagnostics before removing any of them.
\item \textbf{Which seed constraints survive translation?} For C, identify
the actual controller constraints that are invariant under translating all
coordinates. An invariant constraint could be imposed on a seed before
solving its sum. A noninvariant safe-position constraint still needs a
separate construction.
\item \textbf{What is the right normalization statement for higher lifts?}
Generalize D to a known finite value \(v_p(g-1)=s\), maintaining the exact
natural-subtraction and zero-exponent behavior. Assess whether the current
auxiliary-order clients actually need that generality before introducing
another interface.
\end{enumerate}

\section*{Contents of the packet}

The packet contains this LaTeX manuscript and PDF, checked Lean prototypes,
the successfully replayed prime-power production experiment, portable
replay instructions, a source manifest, a novelty ledger, the prior-packet
inventory, validation logs, and a SHA-256 manifest. It deliberately omits
large source snapshots, downloaded runtimes, and compiled Mathlib caches.
The requested source pin and dependency revisions are sufficient to
reconstruct the validation environment.

\begin{multicols}{2}
\begin{thebibliography}{99}
\footnotesize
\bibitem{prior} ProveIt, \emph{Lean proof-simplification proposals}, pinned
packet index and adoption record.
\href{https://github.com/VladimirReshetnikov/ProveIt/tree/e1c8914bc7064d9a8e7c06171a18069a949badd0/Computability/DiophantineConsistency/research/lean-proof-simplification}{Open pinned source}.
\bibitem{incoming} ProveIt, \emph{Incoming reports}, pinned archive queue.
\href{https://github.com/VladimirReshetnikov/ProveIt/tree/e1c8914bc7064d9a8e7c06171a18069a949badd0/docs/incoming}{Open pinned source}.
\bibitem{lateprior} ProveIt, late research index at
\code{ae736b383e59c32c3c3f81910ef1c1197cbfbd44}.
\href{https://github.com/VladimirReshetnikov/ProveIt/tree/ae736b383e59c32c3c3f81910ef1c1197cbfbd44/Computability/DiophantineConsistency/research/lean-proof-simplification}{Open late pinned index}.
\bibitem{lateincoming} ProveIt, incoming directory at the late cutoff.
\href{https://github.com/VladimirReshetnikov/ProveIt/tree/ae736b383e59c32c3c3f81910ef1c1197cbfbd44/docs/incoming}{Open late pinned directory}.
\bibitem{updates} ProveIt, packet \code{2026-10-26zd}, especially
finite prime-factor order certificates and conditional numeric control-slot invariants.
\href{https://github.com/VladimirReshetnikov/ProveIt/blob/ae736b383e59c32c3c3f81910ef1c1197cbfbd44/Computability/DiophantineConsistency/research/lean-proof-simplification/2026-10-26zd/report.tex}{Open late pinned report}.
\bibitem{terminal} ProveIt, packet \code{2026-10-26zh}, exact terminal-band
convolution response.
\href{https://github.com/VladimirReshetnikov/ProveIt/blob/ae736b383e59c32c3c3f81910ef1c1197cbfbd44/Computability/DiophantineConsistency/research/lean-proof-simplification/2026-10-26zh/report/06-word.tex}{Open late pinned section}.
\bibitem{mersenne} ProveIt, packet \code{2026-10-26zg}, Mersenne
exponent classification and related local arithmetic.
\href{https://github.com/VladimirReshetnikov/ProveIt/blob/ae736b383e59c32c3c3f81910ef1c1197cbfbd44/Computability/DiophantineConsistency/research/lean-proof-simplification/2026-10-26zg/report/sections/03-local.tex}{Open late pinned section}.
\bibitem{safety} ProveIt, \code{CurrentDummySafety2011.lean} at the late cutoff.
\href{https://github.com/VladimirReshetnikov/ProveIt/blob/ae736b383e59c32c3c3f81910ef1c1197cbfbd44/Computability/DiophantineConsistency/Lean/DiophantineConsistency/CurrentDummySafety2011.lean}{Open late pinned source}.
\bibitem{reflection} ProveIt, \code{CurrentReflectedGeometry2011.lean},
especially \code{reflected_source} and \code{reflected_cell}.
\href{https://github.com/VladimirReshetnikov/ProveIt/blob/e1c8914bc7064d9a8e7c06171a18069a949badd0/Computability/DiophantineConsistency/Lean/DiophantineConsistency/CurrentReflectedGeometry2011.lean}{Open pinned source}.
\bibitem{topbounds} ProveIt, \code{CurrentWordTopBounds2011.lean}.
\href{https://github.com/VladimirReshetnikov/ProveIt/blob/e1c8914bc7064d9a8e7c06171a18069a949badd0/Computability/DiophantineConsistency/Lean/DiophantineConsistency/CurrentWordTopBounds2011.lean}{Open pinned source}.
\bibitem{units} ProveIt, \code{CurrentControllerUnits2011.lean}.
\href{https://github.com/VladimirReshetnikov/ProveIt/blob/e1c8914bc7064d9a8e7c06171a18069a949badd0/Computability/DiophantineConsistency/Lean/DiophantineConsistency/CurrentControllerUnits2011.lean}{Open pinned source}.
\bibitem{intervals} ProveIt, \code{ControllerIntervalSums.lean}.
\href{https://github.com/VladimirReshetnikov/ProveIt/blob/e1c8914bc7064d9a8e7c06171a18069a949badd0/Computability/DiophantineConsistency/Lean/DiophantineConsistency/ControllerIntervalSums.lean}{Open pinned source}.
\bibitem{orbit} ProveIt, \code{ControllerOrbitTransport.lean}.
\href{https://github.com/VladimirReshetnikov/ProveIt/blob/e1c8914bc7064d9a8e7c06171a18069a949badd0/Computability/DiophantineConsistency/Lean/DiophantineConsistency/ControllerOrbitTransport.lean}{Open pinned source}.
\bibitem{order} ProveIt, \code{PrimePowerOrder.lean}.
\href{https://github.com/VladimirReshetnikov/ProveIt/blob/e1c8914bc7064d9a8e7c06171a18069a949badd0/Computability/DiophantineConsistency/Lean/DiophantineConsistency/PrimePowerOrder.lean}{Open pinned source}.
\end{thebibliography}
\end{multicols}
\end{document}
